\documentclass[11pt,letterpaper]{article}
\usepackage[margin=1in]{geometry}
\usepackage{amsmath,amssymb,bm}
\usepackage{booktabs}
\usepackage{enumitem}
\usepackage{xcolor}
\usepackage[colorlinks=true,linkcolor=blue!70!black,urlcolor=blue!70!black]{hyperref}
\usepackage{fancyhdr}
\usepackage{titlesec}
\usepackage{mathtools}
\usepackage{tcolorbox}

\pagestyle{fancy}
\fancyhf{}
\rhead{Building a BFGS Optimizer --- Guided Exercises}
\lhead{}
\rfoot{Page \thepage}
\setlength{\headheight}{14pt}

\setlength{\parskip}{5pt}
\setlength{\parindent}{0pt}

\titlespacing*{\section}{0pt}{14pt}{6pt}
\titlespacing*{\subsection}{0pt}{10pt}{4pt}
\titlespacing*{\subsubsection}{0pt}{8pt}{3pt}

\newcommand{\R}{\mathbb{R}}
\newcommand{\grad}{\nabla}
\newcommand{\T}{^{\mathsf{T}}}
\newcommand{\half}{\tfrac12}

% Exercise environment
\newtcolorbox{exercise}[1]{
  colback=blue!5!white, colframe=blue!50!black, fonttitle=\bfseries,
  title={Exercise: #1}, sharp corners, boxrule=0.8pt,
  left=6pt, right=6pt, top=4pt, bottom=4pt}

% Hint environment
\newtcolorbox{hint}{
  colback=yellow!5!white, colframe=yellow!60!black, fonttitle=\bfseries,
  title={Hint}, sharp corners, boxrule=0.5pt,
  left=6pt, right=6pt, top=4pt, bottom=4pt}

% Checkpoint environment
\newtcolorbox{checkpoint}[1]{
  colback=green!5!white, colframe=green!50!black, fonttitle=\bfseries,
  title={Checkpoint: #1}, sharp corners, boxrule=0.8pt,
  left=6pt, right=6pt, top=4pt, bottom=4pt}

\title{\textbf{Building a BFGS Optimizer}\\[8pt]
\large Guided Coding Exercises\\[4pt]
\normalsize Inverse-Hessian Quasi-Newton on a Convex Quadratic}
\author{}
\date{June 2026}

\begin{document}
\maketitle
\thispagestyle{fancy}

\tableofcontents
\newpage

%======================================================================
\section{Introduction}
%======================================================================

This document walks you through implementing the \textbf{BFGS} quasi-Newton method
from scratch in MATLAB, one piece at a time, and using it to minimize a convex
quadratic. By the end you will have a small, reusable optimizer whose search step is
a single matrix--vector product, and you will have seen the property that makes BFGS
on a quadratic so clean: \emph{finite termination in at most $n$ steps} with an exact
line search.

We minimize \textbf{Problem 1} of Kanamori \& Ohara (arXiv:1010.2846, \S6):
\begin{equation}\label{eq:prob1}
\min_{x\in\R^n}\ f(x)=\half\,x\T A x - e\T x,
\qquad e=(1,\dots,1)\T,\quad
A=\operatorname{tridiag}(-1,2,-1),
\end{equation}
the second-difference (1-D Laplacian) matrix --- symmetric positive definite, so $f$
is strongly convex with the unique minimizer $x^\ast=A^{-1}e$.

\subsection{Prerequisites}
Before starting you should have:
\begin{itemize}[nosep]
  \item The theory note \texttt{newton\_quasinewton\_bfgs.tex} read --- especially the
    secant equation, the curvature condition $s\T y>0$, and the BFGS inverse-Hessian
    update (the ``$V H V\T+\rho ss\T$'' form).
  \item Basic MATLAB: functions, anonymous-function handles, \texttt{struct}s.
\end{itemize}

\subsection{What You Will Build}
\begin{center}
\begin{tabular}{lp{8.4cm}}
\toprule
\textbf{File} & \textbf{Purpose} \\
\midrule
\texttt{prob1\_obj.m}     & Objective $f$ and gradient $\grad f$ (Problem 1) \\
\texttt{bfgs\_update\_H.m} & One BFGS update of the inverse-Hessian model $H$ \\
\texttt{qn\_minimize\_H.m} & The quasi-Newton driver loop (step, line search, update) \\
\texttt{demo\_prob1.m}    & Build $A,e$, run the optimizer, report convergence \\
\texttt{dfp\_update\_H.m} & \emph{(optional, Exercise 6)} the DFP update, for comparison \\
\bottomrule
\end{tabular}
\end{center}

Work in a scratch folder so your build stays separate from the reference solutions:
create a \texttt{mytry/} subfolder and write your files there. Reference
implementations live in the parent \texttt{quasiNewton\_matlab/} folder --- consult
them only \emph{after} attempting each exercise.

\subsection{Notation Recap}
The quasi-Newton iteration maintains an approximation $H_k$ to the \emph{inverse}
Hessian $(\grad^2 f)^{-1}$. At each step,
\begin{align*}
p_k &= -H_k\,\grad f(x_k) && \text{(search direction --- a matrix--vector product, no solve)}\\
x_{k+1} &= x_k + \alpha_k\,p_k && \text{(line search for the step length $\alpha_k$)}\\
H_{k+1} &= \textsf{BFGS}(H_k,\,s_k,\,y_k) && s_k=x_{k+1}-x_k,\ \ y_k=\grad f(x_{k+1})-\grad f(x_k).
\end{align*}
The pair $(s_k,y_k)$ feeds the \textbf{secant equation} $H_{k+1}y_k=s_k$; the
\textbf{curvature condition} $s_k\T y_k>0$ is what lets the update stay symmetric
positive definite (and hence keep $p_k$ a descent direction).

%======================================================================
\section{Phase A: The Objective}
%======================================================================

%----------------------------------------------------------------------
\subsection{Exercise 1: Objective and Gradient}
%----------------------------------------------------------------------

The optimizer needs a function returning $f(x)$ and, on demand, $\grad f(x)$. For
\eqref{eq:prob1}, $\grad f(x)=Ax-e$ and $\grad^2 f=A$.

\begin{exercise}{1. Write prob1\_obj.m}
Create \texttt{prob1\_obj.m}:
\begin{verbatim}
  function [f, g] = prob1_obj(x, A, e)
\end{verbatim}
Return $f=\half x\T A x - e\T x$, and compute $g=Ax-e$ \emph{only when requested}.
\end{exercise}

\begin{hint}
Guard the gradient with \texttt{nargout} so the same handle serves both a
value-only line-search call and a value-and-gradient step call:
\begin{verbatim}
  f = 0.5 * (x' * A * x) - e' * x;
  if nargout > 1
      g = A * x - e;
  end
\end{verbatim}
This matters: if you build \texttt{f}/\texttt{g} with \texttt{deal(...)} inside an
anonymous function, a single-output call (\texttt{ft = fun(xt)}) will error, because
\texttt{deal} requires the output count to match. A proper function file with an
\texttt{if nargout} guard avoids this.
\end{hint}

\begin{checkpoint}{Gradient matches finite differences}
Build a small instance and check the analytic gradient against central differences:
\begin{verbatim}
  n = 6;  A = full(spdiags(ones(n,1)*[-1 2 -1], -1:1, n, n));  e = ones(n,1);
  x = randn(n,1);  [~, g] = prob1_obj(x, A, e);
  gfd = zeros(n,1);  h = 1e-6;
  for i = 1:n
      d = zeros(n,1); d(i) = h;
      gfd(i) = (prob1_obj(x+d,A,e) - prob1_obj(x-d,A,e)) / (2*h);
  end
  fprintf('grad error: %.3e\n', norm(g - gfd));   % expect < 1e-7
\end{verbatim}
Also confirm the minimizer: \texttt{xstar = A\textbackslash e;} then
\texttt{[\~{},gs] = prob1\_obj(xstar,A,e); norm(gs)} should be $\sim 10^{-14}$
(the gradient vanishes at $x^\ast$).
\end{checkpoint}

%======================================================================
\section{Phase B: The BFGS Update}
%======================================================================

%----------------------------------------------------------------------
\subsection{Exercise 2: The Inverse-Hessian Update}
%----------------------------------------------------------------------

This is the heart of the method. The BFGS update of the inverse Hessian is
\begin{equation}\label{eq:bfgsH}
H_{k+1}=V H_k V\T+\rho\,s s\T,
\qquad V=I-\rho\,s y\T,\qquad \rho=\frac1{s\T y}.
\end{equation}
Two properties fall out of this form for free, and they are the whole reason to write
it this way:
\begin{itemize}[nosep]
  \item \textbf{Symmetry} --- it is a congruence $V H_k V\T$ (preserves symmetry)
    plus the symmetric rank-one $\rho ss\T$.
  \item \textbf{Positive definiteness, when $s\T y>0$} --- for any $z\neq0$,
    $z\T H_{k+1}z=(V\T z)\T H_k(V\T z)+\rho(s\T z)^2$; both terms are $\ge0$ and
    cannot vanish together (if $s\T z=0$ then $V\T z=z$, making the first term
    $z\T H_k z>0$).
\end{itemize}

\begin{exercise}{2. Write bfgs\_update\_H.m}
Create \texttt{bfgs\_update\_H.m}:
\begin{verbatim}
  function H1 = bfgs_update_H(H, s, y)
\end{verbatim}
Steps:
\begin{enumerate}[nosep]
  \item Compute \texttt{sy = s' * y}. If \texttt{sy <= 0}, the curvature condition
    fails and no PD update exists --- \texttt{error(...)} rather than silently return
    a non-PD matrix. (The caller is responsible for guaranteeing $s\T y>0$, or for
    skipping the update; see Exercise~4.)
  \item Form \texttt{rho = 1/sy} and \texttt{V = eye(n) - rho*(s*y')}.
  \item Return \texttt{H1 = V*H*V' + rho*(s*s')}.
  \item Symmetrize against round-off: \texttt{H1 = (H1 + H1')/2}.
\end{enumerate}
\end{exercise}

\begin{hint}
Get $n$ from the step vector, \texttt{n = numel(s)}. The whole body is five lines.
Do not try to be clever and expand $VHV\T$ by hand --- the matrix form is clearer and
just as fast at this scale.
\end{hint}

\begin{checkpoint}{The update satisfies its three requirements}
Pick a random SPD \texttt{H} and random \texttt{s}, \texttt{y} with \texttt{s'*y>0},
then verify:
\begin{verbatim}
  n = 5;  M = randn(n); H = M*M' + n*eye(n);
  s = randn(n,1); y = randn(n,1); if s'*y <= 0, y = -y; end
  H1 = bfgs_update_H(H, s, y);
  fprintf('secant   ||H1*y - s|| = %.2e\n', norm(H1*y - s));   % ~1e-15
  fprintf('symmetry ||H1 - H1''|| = %.2e\n', norm(H1 - H1'));  % ~1e-16
  fprintf('min eig(H1) = %.4f  (PD if > 0)\n', min(eig(H1)));
\end{verbatim}
All three should pass: secant residual at machine precision, symmetric, and
positive definite.
\end{checkpoint}

%======================================================================
\section{Phase C: The Line Search}
%======================================================================

%----------------------------------------------------------------------
\subsection{Exercise 3: Exact Line Search for the Quadratic}
%----------------------------------------------------------------------

Given a descent direction $p$, the line search picks the step length $\alpha$. For a
quadratic, the minimizer along $p$ has a closed form: differentiate
$\phi(\alpha)=f(x+\alpha p)$ and set $\phi'(\alpha)=0$,
\begin{equation}\label{eq:exactls}
\phi'(\alpha)=\grad f(x+\alpha p)\T p=(g+\alpha A p)\T p=0
\quad\Longrightarrow\quad
\boxed{\;\alpha^\ast=-\frac{g\T p}{p\T A p}\;}.
\end{equation}
This is exact (no iteration), always positive for a descent direction with $A\succ0$,
and it is what makes the finite-termination property in Phase~E appear.

\begin{exercise}{3. Build the exact line-search handle}
You will pass the line search to the driver as a function handle with signature
\texttt{alpha = ls(fun, x, f, g, p)}. For the quadratic, write it inline in the demo
(Exercise~5) using \eqref{eq:exactls}:
\begin{verbatim}
  exact_ls = @(fun, x, f, g, p) -(g' * p) / (p' * A * p);
\end{verbatim}
(The general, non-quadratic case needs a backtracking or Wolfe search instead; that
is the ``what comes next'' of \S\ref{sec:next}. Keep the driver agnostic by always
taking a line-search handle.)
\end{exercise}

\begin{checkpoint}{Exact step zeroes the directional derivative}
For a random $x$ and descent direction $p=-g$:
\begin{verbatim}
  [f,g] = prob1_obj(x,A,e);  p = -g;
  a = -(g'*p)/(p'*A*p);
  [~,gnew] = prob1_obj(x + a*p, A, e);
  fprintf('phi''(a*) = %.2e\n', gnew'*p);   % ~1e-14: slope is zero at the min
\end{verbatim}
The directional derivative at the chosen step should be (numerically) zero --- you
landed at the bottom of the 1-D slice.
\end{checkpoint}

%======================================================================
\section{Phase D: The Driver Loop}
%======================================================================

%----------------------------------------------------------------------
\subsection{Exercise 4: The Quasi-Newton Loop}
%----------------------------------------------------------------------

Now assemble the iteration. The driver is \emph{generic} --- it takes the objective,
a start point, an update handle (so you can later swap BFGS for DFP), and options
(including the line-search handle).

\begin{exercise}{4. Write qn\_minimize\_H.m}
Create \texttt{qn\_minimize\_H.m}:
\begin{verbatim}
  function [x, info] = qn_minimize_H(fun, x0, update_fn, opts)
\end{verbatim}
Steps:
\begin{enumerate}[nosep]
  \item Read options with defaults: \texttt{H0} (default \texttt{eye(n)}),
    \texttt{tol} (\texttt{1e-8}), \texttt{maxit} (\texttt{1000}), and the
    \texttt{linesearch} handle (default: a backtracking-Armijo helper).
  \item Initialize \texttt{x = x0(:)}, \texttt{[f,g] = fun(x)}, \texttt{H = H0}.
  \item Loop while \texttt{k < maxit}:
    \begin{itemize}[nosep]
      \item If \texttt{norm(g) <= tol}, mark converged and break.
      \item Direction \texttt{p = -H * g}. Guard: if \texttt{g'*p >= 0} (not a descent
        direction --- should not happen with $H\succ0$, but round-off and bad data can
        cause it), reset \texttt{p = -g; H = eye(n)}.
      \item \texttt{alpha = ls(fun, x, f, g, p)}; \texttt{xnew = x + alpha*p}.
      \item \texttt{[fnew,gnew] = fun(xnew)}; \texttt{s = xnew-x}; \texttt{y = gnew-g}.
      \item \textbf{Curvature guard:} update only if
        \texttt{s'*y > 1e-12*norm(s)*norm(y)}; else increment a \texttt{skips} counter
        and leave $H$ unchanged (keeps $H\succ0$).
      \item Advance: \texttt{x=xnew; f=fnew; g=gnew; k=k+1}; record \texttt{norm(g)}.
    \end{itemize}
  \item Return \texttt{x} and an \texttt{info} struct: \texttt{iters},
    \texttt{gnorm}, \texttt{gnorm\_hist}, \texttt{converged}, \texttt{skips}.
\end{enumerate}
\end{exercise}

\begin{hint}
Default-option pattern and the curvature guard:
\begin{verbatim}
  if nargin < 4, opts = struct(); end
  n   = numel(x0);
  H0  = getfield_default(opts,'H0',eye(n));
  tol = getfield_default(opts,'tol',1e-8);
  maxit = getfield_default(opts,'maxit',1000);
  c1  = getfield_default(opts,'c1',1e-4);   % Armijo constant
  bt  = getfield_default(opts,'bt',0.5);    % backtracking factor
  ls  = getfield_default(opts,'linesearch', ...
            @(f_,x_,f0_,g_,p_) armijo(f_,x_,f0_,g_,p_,c1,bt));
  ...
  if s'*y > 1e-12 * norm(s) * norm(y)
      H = update_fn(H, s, y);
  else
      skips = skips + 1;
  end
\end{verbatim}
Put two local functions \emph{at the end of the file} (MATLAB requires local
functions after the body). First, the one-line option helper:
\begin{verbatim}
  function v = getfield_default(s, name, dflt)
      if isfield(s,name) && ~isempty(s.(name)), v = s.(name); else, v = dflt; end
  end
\end{verbatim}
Second, the default line search --- \textbf{backtracking Armijo}. It enforces the
\emph{sufficient-decrease} (Armijo) condition $f(x+\alpha p)\le f(x)+c_1\,\alpha\,g\T p$
(with $g\T p<0$ for a descent direction): start at $\alpha=1$ and shrink by $bt$ until
it holds.
\begin{verbatim}
  function alpha = armijo(fun, x, f0, g, p, c1, bt)
      alpha = 1;  gtp = g' * p;
      while true
          ft = fun(x + alpha*p);
          if ft <= f0 + c1*alpha*gtp || alpha < 1e-16, break; end
          alpha = bt * alpha;
      end
  end
\end{verbatim}
This default makes the driver usable on a general function; for the quadratic of this
tutorial you will override it with the \emph{exact} line search (Exercise~3, plugged in
at Exercise~5), which is what produces the $n$-step termination. Taking the update as a
handle (\texttt{update\_fn}) is what lets you reuse this same loop for DFP later.
\end{hint}

\begin{checkpoint}{Loop runs and the direction is a descent direction}
Before worrying about convergence speed, confirm the loop runs one BFGS step and that
the very first direction (with $H_0=I$, so $p_0=-g_0$) decreases $f$:
\begin{verbatim}
  fun = @(x) prob1_obj(x, A, e);
  [f0,g0] = fun(x0);  p0 = -g0;
  fprintf('g0''*p0 = %.3e  (must be < 0)\n', g0'*p0);
\end{verbatim}
A negative slope confirms descent. If you ran with a backtracking line search and saw
many \texttt{skips} on this quadratic, that is expected --- Armijo alone forfeits the
finite-termination property. Phase~E fixes that with the exact line search.
\end{checkpoint}

%======================================================================
\section{Phase E: Run It, and the $n$-Step Surprise}\label{sec:phaseE}
%======================================================================

%----------------------------------------------------------------------
\subsection{Exercise 5: The Demo Driver}
%----------------------------------------------------------------------

\begin{exercise}{5. Write demo\_prob1.m}
Create \texttt{demo\_prob1.m}:
\begin{enumerate}[nosep]
  \item Set \texttt{rng(0)} for reproducibility; choose \texttt{n = 20}.
  \item Build \texttt{A = full(spdiags(ones(n,1)*[-1 2 -1], -1:1, n, n))} and
    \texttt{e = ones(n,1)}; set \texttt{fun = @(x) prob1\_obj(x, A, e)}.
  \item Compute \texttt{xstar = A\textbackslash e} and a start point
    \texttt{x0 = sqrt(10)*randn(n,1)}.
  \item Build the exact line search \eqref{eq:exactls} and put it in \texttt{opts}:
\begin{verbatim}
  exact_ls = @(fun,x,f,g,p) -(g'*p)/(p'*A*p);
  opts = struct('tol', n*1e-10, 'maxit', 2000, 'linesearch', exact_ls);
\end{verbatim}
  \item Run \texttt{[xB,infoB] = qn\_minimize\_H(fun, x0, @bfgs\_update\_H, opts)} and
    print \texttt{infoB.iters}, \texttt{infoB.gnorm}, \texttt{norm(xB-xstar)}, and
    \texttt{infoB.skips}.
\end{enumerate}
\end{exercise}

\begin{checkpoint}{Finite termination: BFGS converges in $\le n$ steps}
With the exact line search, $n=20$, and \texttt{rng(0)} you should see:
\begin{verbatim}
  method   iters     final |g|        ||x-x*||   skips
  BFGS        20     ~3e-14          ~6e-14          0
\end{verbatim}
\textbf{Exactly $n=20$ iterations}, gradient driven to $\sim10^{-14}$, zero skips.
This is the conjugate-direction finite-termination property: BFGS with an exact line
search minimizes a strongly convex quadratic in at most $n$ steps. If you instead see
hundreds of iterations, you are almost certainly \emph{not} using the exact line
search (a backtracking Armijo search loses this property --- see Exercise~6).
\end{checkpoint}

%----------------------------------------------------------------------
\subsection{Exercise 6 (stretch): DFP, and why the line search matters}
%----------------------------------------------------------------------

\begin{exercise}{6. Swap in DFP and compare}
Write \texttt{dfp\_update\_H.m} --- the DFP inverse-Hessian update,
\[
  H_{k+1}=H_k+\frac{s s\T}{s\T y}-\frac{H_k y\,y\T H_k}{y\T H_k y},
\]
with the same \texttt{s'*y>0} guard. Run the demo with \texttt{@dfp\_update\_H} as
well, under \emph{two} line searches:
\begin{enumerate}[nosep]
  \item the exact line search of \eqref{eq:exactls};
  \item a backtracking-Armijo search (the driver's default).
\end{enumerate}
Compare iteration counts for BFGS vs DFP in each case.
\end{exercise}

\begin{checkpoint}{The robustness gap}
With the \emph{exact} line search, both BFGS and DFP terminate in $n=20$ steps (the
finite-termination property holds for the whole Broyden family). With the
\emph{inexact} backtracking search, BFGS still converges quickly and stably while DFP
slows dramatically (many iterations, many curvature skips). This \emph{qualitatively
mirrors} the Kanamori \& Ohara result (their Table~3 sweeps line-search \emph{noise}
$h$ with a near-exact \texttt{fminbnd} search rather than this plain Armijo, so it is
a foreshadowing, not a reproduction of their exact protocol): BFGS is robust to an
inexact line search; DFP is not. It is the empirical face of the variational/asymmetry
story in the theory note --- and the reason BFGS, not DFP, is the default quasi-Newton
method. (To reproduce their experiment proper, inject the $\tilde s=(1+\varepsilon)s$,
$\varepsilon\sim\mathcal U[-h,h]$ perturbation and sweep $h$; see
\S\ref{sec:next}.)
\end{checkpoint}

%======================================================================
\section{What Comes Next}\label{sec:next}
%======================================================================

You now have a working BFGS optimizer whose step is a matrix--vector product. Natural
extensions, in roughly increasing difficulty:

\begin{enumerate}
  \item \textbf{A real line search.} Replace the quadratic-only exact step with a
    backtracking line search satisfying the \emph{Wolfe} conditions. The curvature
    (second Wolfe) condition is precisely what guarantees $s\T y>0$ on a general
    function --- the hypothesis our update needs.
  \item \textbf{Problem 2 (nonlinear, non-convex).} Add the trigonometric
    boundary-value objective $f(x)=\half x\T Ax-e\T x-\frac1{(n+1)^2}\sum_i(2x_i+\cos x_i)$.
    Its Hessian is $A+\frac1{(n+1)^2}\operatorname{diag}(\cos x_i)$; the secant
    equation is now a genuine model, not an identity. A Wolfe line search becomes
    necessary.
  \item \textbf{The robustness experiment.} Reproduce Kanamori \& Ohara Table~3:
    perturb the step $\tilde s_k=(1+\varepsilon)s_k$ with $\varepsilon\sim
    \mathcal U[-h,h]$, sweep $h\in[0,0.3]$, and plot iteration count vs $h$ for BFGS
    and DFP on both problems.
  \item \textbf{L-BFGS.} For large $n$, never form $H$ at all: store the last $m$
    pairs $(s_k,y_k)$ and apply $H_k$ implicitly via the two-loop recursion. This
    removes the $O(n^2)$ storage and is what makes quasi-Newton scale to
    $n=10^6$.
\end{enumerate}

Each builds on the loop you just wrote without restructuring its core: the direction
$p=-Hg$, the line search, the curvature-guarded update.

\end{document}
